%%% BEGIN fig:method-sudakov-plane
\begin{figure}[htbp]
\centering
\begin{tikzpicture}[fd,
  fb lab/.style={font=\footnotesize,inner sep=1.5pt},
  fb tick/.style={font=\scriptsize},
  fb grid/.style={draw=gray!45,line width=0.4pt,dash pattern=on 1pt off 2pt},
  fb hyp/.style={draw=black!60,line width=0.6pt},
  fb patch/.style={circle,draw=black!70,line width=0.6pt,minimum size=0.62cm,inner sep=0pt,font=\small},
  fb pole/.pic={\draw[line width=1pt,line cap=round] (-0.085,-0.085)--(0.085,0.085) (-0.085,0.085)--(0.085,-0.085);},
  fb contour/.style={draw=blue!60!black,line width=1.0pt},
  fb carrow/.style={postaction={decorate,decoration={markings,
      mark=at position #1 with {\arrow{Stealth[length=5.5pt,width=4.5pt]}}}}}]
  % ======================= (a) the (ln k^+, ln k^-) plane =======================
  \begin{scope}[x=1.3cm,y=1.3cm]
    % soft--collinear overlap strips (hatched)
    \fill[pattern=north east lines,pattern color=green!45!black] (-2,-2.17) rectangle (0,-1.83);
    \fill[pattern=north east lines,pattern color=green!45!black] (-2.17,-2) rectangle (-1.83,0);
    % axes, ticks, grid
    \draw[lcaxis] (-2.8,-2.8) -- (0.8,-2.8) node[below left=1pt and -6pt,font=\small] {$\ln|k^+|$};
    \draw[lcaxis] (-2.8,-2.8) -- (-2.8,0.85) node[left,font=\small] {$\ln|k^-|$};
    \foreach \v/\t in {-2/{\lambda^2Q},-1/{\lambda Q},0/{Q}} {
      \draw[line width=0.5pt] (\v,-2.8) -- ++(0,0.07);
      \node[fb tick,below] at (\v,-2.8) {$\t$};
      \draw[line width=0.5pt] (-2.8,\v) -- ++(0.07,0);
      \node[fb tick,left] at (-2.8,\v) {$\t$};
      \draw[fb grid] (\v,-2.73) -- (\v,0.65);
      \draw[fb grid] (-2.73,\v) -- (0.65,\v);
    }
    % lines of constant k^+ k^-
    \draw[fb hyp] (0.62,-0.62) -- (-0.62,0.62);
    \draw[fb hyp,line width=1.0pt] (0.5,-2.5) -- (-2.5,0.5);
    \node[fb lab,text=black!65,anchor=south east] at (-0.6,0.6) {$k^+k^-\sim Q^2$};
    \node[fb lab,text=black!80,anchor=north west] at (0.5,-2.48) {$k^+k^-\sim m^2$};
    % the family k=(z, lambda^2/z, lambda) between C1 and C2
    \draw[red!60!black,line width=1.6pt,dash pattern=on 4pt off 2pt] (-1.72,-0.28) -- (-0.28,-1.72);
    % regions of the off-shell triangle
    \node[fb patch,fill=gray!35] (H) at (0,0) {$\mathrm H$};
    \node[fb patch,fill=blue!20] (C1) at (0,-2) {$C_1$};
    \node[fb patch,fill=orange!25] (C2) at (-2,0) {$C_2$};
    \node[fb patch,fill=green!25] (S) at (-2,-2) {$S$};
    \draw[black!70,line width=0.6pt,fill=white] (-1,-1) circle[radius=2.3pt];
    \node[font=\scriptsize,inner sep=1pt,anchor=north east] at (-1.05,-1.05) {$(\lambda,\lambda,\lambda)$};
    \node[fb lab,text=red!55!black,anchor=south west,align=left] at (-0.62,-1.08) {$k=(z,\lambda^2/z,\lambda)$\\$\lambda^2\ll z\ll1$：见 (b)};
    \node[fb lab,text=green!30!black,anchor=north] at (-1,-2.19) {软—共线重叠};
    \node[font=\small,anchor=south west] at (-2.75,0.9) {(a)};
  \end{scope}
  % ======================= (b) complex k^- plane =======================
  \begin{scope}[xshift=7.6cm,yshift=-1.75cm]
    \fill[red!12] (0,0) circle[radius=0.32];
    \draw[lcaxis] (-2.25,0) -- (3.0,0) node[below left=1pt and -4pt,font=\small] {$\Real k^-$};
    \draw[lcaxis] (0,-1.7) -- (0,1.8) node[left,font=\small] {$\Imag k^-$};
    % contour: lower half plane on the left, crossing near k_T^2/(4k^+), upper half plane on the right
    \draw[fb contour,fb carrow=0.16,fb carrow=0.84] (-2.15,-0.9) .. controls (-0.55,-0.9) and (0.45,-0.5) .. (0.95,0)
         .. controls (1.45,0.5) and (2.05,0.9) .. (2.85,0.9);
    \pic at (-0.3,0.2) {fb pole};
    \pic at (1.9,-0.2) {fb pole};
    \fill[blue!60!black] (0.95,0) circle[radius=1.5pt];
    \node[fb lab,anchor=south east] at (-0.36,0.24) {$D_1$};
    \node[fb lab,anchor=north west] at (1.96,-0.24) {$D_g$};
    \draw[gray!50!black,line width=0.4pt] (0.95,0) -- (1.2,0.62) node[above,fb lab] {$\Real k^-\approx\frac{k_T^2}{4k^+}$};
    \node[fb lab,text=red!55!black,anchor=north east,align=right] at (-0.3,-0.12) {$k^-\approx0$：\\软近似失效};
    \node[fb lab,anchor=north,align=left] at (0.4,-1.85) {$D_1$：$k^-\approx-\frac{m^2+k_T^2}{2P}+\ii0$，\quad $D_g$：$k^-=\frac{k_T^2-\ii0}{2k^+}$};
    \node[font=\small,anchor=south west] at (-2.2,1.75) {(b)};
  \end{scope}
\end{tikzpicture}
\caption{单圈离壳三角图（$p_i^2=-m^2$，$\lambda=m/Q$）的区域与 Collins 减除中的软—共线重叠（\cref{ex:method-sudakov-subtraction}）。
(a)~$(\ln|k^+|,\ln|k^-|)$ 平面（以 $\ln(1/\lambda)$ 为单位，$k_T$ 取 $k_T^2\sim|k^+k^-|$ 的典型值）：区域 $\mathrm H$、$C_1$、$C_2$、$S$ 位于正方形的四角，其中 $S\sim Q(\lambda^2,\lambda^2,\lambda^2)$ 即\cref{def:ir-scalings}的超软标度；
粗斜线 $k^+k^-\sim m^2$ 是 $k_T\sim\lambda Q$ 的在壳线 $k_T^2=2k^+k^-$，经过 $C_1$、$C_2$ 与“软中心”$(\lambda,\lambda,\lambda)$；阴影带为软—共线重叠，减除项在重叠区相消的方式由分组表达式决定。
(b)~对 $k=(z,\lambda^2/z,\lambda)$（$\lambda^2\ll z\ll1$，红色虚线），胶子线 $D_g$ 与夸克线 $D_1$ 的 $k^-$ 极点分居实轴两侧、相距 $\sim\lambda^2/z$，围道只能在两极点之间穿过实轴；
软近似只在 $\Real k^-\sim k_T^2/(4k^+)$ 附近穿过实轴的围道上成立，在 $k^-\approx0$ 处失效（\cref{lem:method-i0-compat}）。}
\label{fig:method-sudakov-plane}
\end{figure}
%%% END fig:method-sudakov-plane
